\section{Choosing the Final Strategy}
\label{sec:choice}

The comparison in Section~\ref{ch:results} produced no strategy that was best on every measure and in
every period. The final choice was therefore made against a small set of requirements for a tool that
an ordinary investor would use every month.

\subsection{Requirements}

\begin{enumerate}
  \item \textbf{Good balance of return and risk on unseen data.} Judged by the utility $U$ and the
        worst fall on 2020--2025, not by the period the rule was designed on.
  \item \textbf{Robust to estimation error.} It should not rely heavily on past average returns,
        which proved unreliable (Section~\ref{sec:res-utility}).
  \item \textbf{Personal.} The investor's tolerance for risk should change the advice; otherwise
        there is no reason to ask for it.
  \item \textbf{Adaptive but stable.} It should use the latest data, without swinging after one
        unusual month.
  \item \textbf{Simple to follow.} A monthly instruction of a few purchases, with selling only when
        the portfolio has clearly drifted.
\end{enumerate}

\subsection{Why the other strategies were not chosen}

\begin{table}[H]
\centering
\caption{The candidates against the requirements}
\label{tab:choice}
\small
\begin{tabularx}{\textwidth}{@{}l>{\raggedright\arraybackslash}X>{\raggedright\arraybackslash}X@{}}
\toprule
Strategy & Strength & Why it was not the final choice \\
\midrule
1 Equity (100\% Nifty) & Highest return in 2010--2019 & Largest falls ($-23.7\%$, $-28.2\%$); never had the highest $U$ for any investor \\
2 Equal-weight & Best return per unit of risk out of sample; very simple & Ignores the data and the investor: the same \nicefrac13 split for everyone, always \\
3 60/20/20 & Best in sample (2010--2019) & Fell to 4th of eight by $U$ out of sample for every profile; a fixed mix that suited one decade \\
4--6 Optimisers vs 0\% & Tiny falls ($-1.7\%$ to $-3.2\%$) & About 70\% in cash: earned little more than the liquid fund itself \\
6b Fixed Optimized & Highest $U$ for every profile out of sample (XIRR 20.3\%) & One 10\% cushion for everyone, so a cautious investor still faces a $-12.2\%$ fall \\
7 Max-Utility & Uses the investor's $\lambda$ directly & Relies on average returns; never ranked first, a clear case of estimation error \\
\bottomrule
\end{tabularx}
\end{table}

Two lessons shaped the final rule. First, the split between Nifty and gold produced by Strategy 6b
(half risk parity, half maximum Sharpe measured above the liquid fund) was the strongest performer out
of sample and the most robust optimiser. Second, the only part of 6b that ignored the investor was the
fixed 10\% liquid share.

\subsection{The final choice: the Personal SIP}

The Personal SIP keeps everything that worked in Strategy 6b and changes only the part that ignored
the investor: the liquid share becomes a cushion chosen by $U$ for each risk level. This follows
Tobin's separation \citep{tobin1958}: how much to keep safe is a personal decision, how to split the
rest is not.

\begin{table}[H]
\centering
\caption{Each design choice of the Personal SIP and the reason for it}
\label{tab:design-choices}
\small
\begin{tabularx}{\textwidth}{@{}>{\raggedright\arraybackslash}p{3.3cm}>{\raggedright\arraybackslash}p{3.3cm}>{\raggedright\arraybackslash}X@{}}
\toprule
Choice & Alternative considered & Reason \\
\midrule
Liquid cushion 10/30/70\% by risk level & One cushion for all (6b) & Makes the advice personal; chosen by $U$ on 2010--2019 only \\
Returns measured above the liquid fund & Above 0\% (Strategies 4--6) & Against 0\% the optimisers simply hold cash \\
Average of risk parity and max Sharpe & Max Sharpe or max utility alone & Risk parity needs no average returns, so it offsets the estimation error of max Sharpe \\
Nifty and gold each 10--70\% & No limits & Prevents an all-in bet on one asset from noisy estimates \\
120-month window, re-fitted monthly & Yearly re-fit & Same performance, and each month's plan uses the latest data \\
Fill gaps first, reset only beyond 5\% & Pro-rata buying; calendar resets & Fewer sales and costs while staying close to target \\
\bottomrule
\end{tabularx}
\end{table}

\paragraph{What the choice costs.} On 2020--2025 the Personal SIP's $U$ was close to Strategy 6b's
for the aggressive investor (9.99\% against 10.25\%), whose cushion is the same 10\%, but lower for
the moderate and conservative investors (Table~\ref{tab:personal}), because their larger cushions
missed part of the gold boom. In return, the conservative investor's worst fall was $-3.4\%$ instead
of $-12.2\%$. The rule therefore does what a cautious investor asks of it: it gives up some return in
strong markets for much smaller losses. Since no rule won in every period, a rule that matches the
investor's own tolerance for risk was preferred to the one that happened to score highest in a single
six-year test.
